\documentclass[11pt]{article}

\usepackage[margin=1in]{geometry}
\usepackage{amsmath,amssymb,amsthm}
\usepackage{algorithm}
\usepackage{algpseudocode}
\usepackage{booktabs}
\usepackage{natbib}
\usepackage[colorlinks=true,linkcolor=blue,citecolor=blue,urlcolor=blue]{hyperref}

\newcommand{\N}{\mathbb{N}}
\newcommand{\R}{\mathbb{R}}
\newcommand{\E}{\mathbb{E}}
\newcommand{\Poi}{\mathrm{Poi}}
\newcommand{\Gam}{\mathrm{Gam}}
\newcommand{\Cat}{\mathrm{Cat}}
\newcommand{\KL}{\mathrm{KL}}

\theoremstyle{definition}
\newtheorem{remark}{Remark}

\title{Mixture-of-Submanifolds Nonnegative Matrix Factorization for Interpretable Sparse Loading Motifs}
\author{TBD}
\date{}

\begin{document}

\maketitle

\begin{abstract}
Nonnegative matrix factorization is widely used to represent a nonnegative data
matrix as low-rank nonnegative factors and loadings. Many sparse variants of
nonnegative matrix factorization encourage individual loading vectors to use
only a small number of factors, but they usually do not model repeated sparsity
patterns as inferential objects. We introduce mixture-of-submanifolds
nonnegative matrix factorization, a probabilistic model for count data that
learns recurring sparse loading motifs. Each observation is assigned, softly, to
a latent motif. Each motif is represented by a small number of
motif dimensions, and each dimension has a posterior distribution over the shared factor
rows. This gives both clustering of observations by sparsity pattern and
posterior uncertainty over the factors defining each motif. The method builds
on a Poisson SuSiE model for one count vector and a matrix-factorization
extension of Poisson SuSiE. We describe the model, a variational empirical-Bayes fitting
procedure, and the simulation studies needed to evaluate factor recovery,
motif recovery, overfitting behavior, and posterior uncertainty.
\end{abstract}

\section{Introduction}

Nonnegative matrix factorization (NMF) seeks a low-rank representation of a
nonnegative matrix $Y \in \N^{N \times M}$ by writing
\[
  Y \approx L F,
\]
where $L \in \R_+^{N \times K}$ contains nonnegative loadings and
$F \in \R_+^{K \times M}$ contains nonnegative factors. In many scientific
applications the rows of $L$ are expected to be sparse: each observation should
be explained by only a few latent factors. Sparse NMF methods address this by
regularizing or constraining the loading matrix. However, sparsity itself is not
the whole inferential target. Often we want to know whether many observations
share the same kind of sparsity pattern, which factors define that pattern, and
whether the evidence distinguishes one possible factor combination from another.

This paper develops MoS-NMF, a mixture-of-submanifolds model for sparse NMF
loadings. The main idea is to treat recurring sparsity patterns as latent
motifs. Each motif represents a small collection of factor
rows that tend to appear together in the loadings. Observations are softly
assigned to motifs, and each motif has posterior uncertainty over which factors
define its dimensions. Thus the model is designed to answer questions such as:
Which sparsity motifs are present? Which observations use each motif? Which
factors define a motif? Are two nearly identical or correlated factors
interchangeable in the posterior?

The statistical challenge is that sparse NMF has several sources of
non-identifiability. Factor rows can be permuted, highly correlated factors can
compete to explain the same counts, and an overfitted factorization can produce
duplicate or redundant factors. A hard sparse loading estimate can hide this
ambiguity by choosing one arbitrary support. We instead represent uncertainty
with posterior probabilities over factor identities, following the spirit of
SuSiE \citep{Wang2020-cr}. The resulting posterior summaries are analogous to
posterior inclusion probabilities, but they operate at the motif level rather
than only at the observation level.

The paper has three methodological layers. Section~\ref{sec:pois-susie}
introduces Poisson SuSiE for one count vector and MF-Poisson-SuSiE for a count
matrix. Section~\ref{sec:mos-nmf} introduces MoS-NMF, which pools sparse loading
patterns across observations through latent motifs. Section~\ref{sec:identifiability}
is reserved for identifiability theory. Section~\ref{sec:experiments} describes
the experiments needed for a JMLR-quality empirical evaluation.

\section{Poisson SuSiE and MF-Poisson-SuSiE}
\label{sec:pois-susie}

\subsection{Single-vector Poisson SuSiE}

We begin with one count vector $y \in \N^M$ and a known nonnegative dictionary
$F \in \R_+^{K \times M}$. Throughout this paper we use row-normalized factors,
\[
  F_{km} \geq 0,
  \qquad
  \sum_{m=1}^M F_{km}=1,
  \qquad k=1,\ldots,K.
\]
Poisson SuSiE represents the Poisson mean as a sparse sum of $D$ loading
dimensions:
\[
  y_m \mid \lambda,\gamma,F
  \sim
  \Poi(\mu_m),
  \qquad
  \mu_m =
  \sum_{d=1}^D \lambda_d F_{\gamma_d m},
  \qquad m=1,\ldots,M.
\]
For each dimension $d$,
\[
  \gamma_d \sim \Cat(1/K,\ldots,1/K),
  \qquad
  \lambda_d \mid \gamma_d \sim \Gam(\alpha_d^0,\beta_d^0),
\]
where the Gamma distribution is parameterized by shape and rate. The latent
index $\gamma_d$ chooses one factor row, and $\lambda_d$ is the nonnegative
loading size for dimension $d$.

The variational approximation factorizes over dimensions:
\[
  q(\lambda,\gamma)
  =
  \prod_{d=1}^D q(\gamma_d)q(\lambda_d \mid \gamma_d),
\]
with
\[
  q(\gamma_d=k)=\bar{\gamma}_{dk},
  \qquad
  q(\lambda_d \mid \gamma_d=k)=\Gam(\alpha_{dk},\beta_{dk}).
\]
We also introduce allocation variables $\xi_{dm}$ satisfying
$\xi_{dm}\geq 0$ and $\sum_{d=1}^D \xi_{dm}=1$ for each coordinate $m$. They
allocate the observed count $y_m$ across the $D$ dimensions. Jensen's inequality
gives the lower bound contribution
\[
  y_m \log\left(\sum_{d=1}^D \lambda_d F_{\gamma_d m}\right)
  \geq
  \sum_{d=1}^D
  \xi_{dm}y_m
  \log\left(\frac{\lambda_d F_{\gamma_d m}}{\xi_{dm}}\right).
\]

The resulting coordinate updates have a simple form. The allocation update is
\[
  \xi_{dm}
  \propto
  \exp\left[
    \sum_{k=1}^K \bar{\gamma}_{dk}
    \{\psi(\alpha_{dk})-\log\beta_{dk}\}
    +
    \sum_{k=1}^K \bar{\gamma}_{dk}\log F_{km}
  \right],
\]
normalized over $d=1,\ldots,D$. The Gamma variational parameters are
\[
  \alpha_{dk}
  \leftarrow
  \alpha_d^0+\sum_{m=1}^M \xi_{dm}y_m,
  \qquad
  \beta_{dk}
  \leftarrow
  \beta_d^0+\sum_{m=1}^M F_{km}.
\]
The posterior probability that dimension $d$ uses factor $k$ is updated by a
softmax over the terms in the lower bound depending on $\gamma_d=k$:
\[
  \bar{\gamma}_{dk}
  \propto
  \exp\left[
    \sum_{m=1}^M \xi_{dm}y_m \log F_{km}
    -
    \left(\alpha_d^0+\sum_{m=1}^M \xi_{dm}y_m\right)
    \log\left(\beta_d^0+\sum_{m=1}^M F_{km}\right)
  \right].
\]
When the rows of $F$ are normalized, the rate term is constant across $k$, but
we keep it visible because it clarifies the role of the Gamma prior and the
scale convention.

\subsection{Empirical-Bayes prior update}

Poisson SuSiE can use empirical Bayes to update the prior mean of each loading
dimension. For dimension $d$, define the posterior mean and posterior mean
log-size
\[
  \sum_{k=1}^K \bar{\gamma}_{dk}\frac{\alpha_{dk}}{\beta_{dk}},
  \qquad
  \sum_{k=1}^K
  \bar{\gamma}_{dk}\{\psi(\alpha_{dk})-\log\beta_{dk}\}.
\]
Maximizing the expected Gamma prior gives the rate update
\[
  \beta_d^0
  =
  \frac{\alpha_d^0}
       {\sum_{k=1}^K \bar{\gamma}_{dk}\alpha_{dk}/\beta_{dk}},
\]
and $\alpha_d^0$ solves
\[
  \psi(\alpha_d^0)
  =
  \sum_{k=1}^K
  \bar{\gamma}_{dk}\{\psi(\alpha_{dk})-\log\beta_{dk}\}
  + \log\beta_d^0.
\]
This empirical-Bayes step is important when $D$ is overspecified: dimensions with
little evidence can have prior mean $\alpha_d^0/\beta_d^0$ close to zero,
which is an automatic relevance determination (ARD) mechanism for selecting the
effective number of dimensions.

\subsection{MF-Poisson-SuSiE}

MF-Poisson-SuSiE applies the same sparse Poisson representation to each row of a
count matrix while learning the shared factor matrix $F$. For observation
$i=1,\ldots,N$,
\[
  Y_{im} \mid \lambda_i,\gamma_i,F
  \sim
  \Poi\left(
    \sum_{d=1}^D \lambda_{id}F_{\gamma_{id}m}
  \right),
  \qquad m=1,\ldots,M.
\]
Each observation has its own sparse support indicators $\gamma_{id}$ and loading
sizes $\lambda_{id}$. The variational posterior uses
\[
  q(\gamma_{id}=k)=\bar{\gamma}_{idk},
  \qquad
  q(\lambda_{id}\mid\gamma_{id}=k)
  =
  \Gam(\alpha_{idk},\beta_{idk}).
\]

For fixed $F$, the updates are the row-wise analogues of the single-vector
updates. When $F$ is unknown, the factor update uses expected allocated counts:
\[
  A_{km}
  =
  \sum_{i=1}^N \sum_{d=1}^D
  \bar{\gamma}_{idk}\xi_{idm}Y_{im}.
\]
With row-normalized factors,
\[
  F_{km}
  \leftarrow
  \frac{A_{km}+\epsilon}
       {\sum_{m'=1}^M(A_{km'}+\epsilon)}.
\]
In practice the implementation initializes $F$ with Poisson NMF, initializes
the sparse posterior from the corresponding loading estimates, and damps the
updates of $\bar{\gamma}$ and $F$. This produces a useful intermediate model:
it learns the factor dictionary and observation-level sparse loading
uncertainty, but it does not yet pool recurring support patterns across
observations. That pooling is the purpose of MoS-NMF.

\section{MoS-NMF}
\label{sec:mos-nmf}

\subsection{Generative model}

MoS-NMF models repeated sparse loading motifs. The observed data are counts
$Y \in \N^{N \times M}$. The factor matrix
$F \in \R_+^{K \times M}$ is row-normalized as above. There are $S$ latent
motifs. We initially fit each motif with a
nominal dimension $D$, but the effective dimension is allowed to be smaller and
can vary by motif.

For motif $s$ and dimension $d$, the factor identity is random:
\[
  \gamma_{sd}
  \sim
  \Cat(\rho_{sd1},\ldots,\rho_{sdK}).
\]
For observation $i$,
\[
  Z_i \sim \Cat(\pi_1,\ldots,\pi_S).
\]
Conditional on $Z_i=s$, the observation has one nonnegative loading size for
each dimension:
\[
  \lambda_{id}\mid Z_i=s
  \sim
  \Gam(\alpha_{sd}^0,\beta_{sd}^0),
  \qquad d=1,\ldots,D.
\]
The observed counts are
\[
  Y_{im}\mid Z_i=s,\lambda_i,\gamma,F
  \sim
  \Poi(\mu_{im}),
  \qquad
  \mu_{im} =
  \sum_{d=1}^D \lambda_{id}F_{\gamma_{sd}m},
  \qquad m=1,\ldots,M.
\]
Equivalently, if $L_i$ denotes the loading vector for observation $i$, then
\[
  L_i
  =
  \sum_{d=1}^D \lambda_{id} e_{\gamma_{Z_i d}},
\]
where $e_k$ is the $k$th coordinate vector in $\R^K$. Thus all observations
assigned to motif $s$ share the same motif-level factor identities
$\gamma_{s1},\ldots,\gamma_{sD}$, but each observation has its own loading
sizes.

\subsection{Variational family}

The variational approximation is
\[
  q(Z,\gamma,\lambda)
  =
  \prod_{i=1}^N q(Z_i)
  \prod_{s=1}^S\prod_{d=1}^D q(\gamma_{sd})
  \prod_{i=1}^N\prod_{s=1}^S\prod_{d=1}^D
  q(\lambda_{id}\mid Z_i=s).
\]
We write
\[
  q(Z_i=s)=\omega_{is},
  \qquad
  q(\gamma_{sd}=k)=\bar{\gamma}_{sdk},
\]
and
\[
  q(\lambda_{id}\mid Z_i=s)
  =
  \Gam(\alpha_{isd},\beta_{sd}).
\]
The rate parameter $\beta_{sd}$ does not depend on $i$ because the Poisson
factor rows are shared and row-normalized; the shape parameter
$\alpha_{isd}$ depends on the counts allocated from observation $i$.

For each candidate motif $s$, we introduce allocation variables
$\xi_{isdm}$ satisfying
\[
  \xi_{isdm}\geq 0,
  \qquad
  \sum_{d=1}^D \xi_{isdm}=1.
\]
They allocate $Y_{im}$ across the dimensions of motif $s$. The implementation
computes $\xi_{isdm}$ in feature blocks rather than storing the full
$N\times S\times D\times M$ tensor.

\subsection{Coordinate updates}

The allocation update is
\[
  \xi_{isdm}
  \propto
  \exp\left[
    \psi(\alpha_{isd})-\log\beta_{sd}
    +
    \sum_{k=1}^K \bar{\gamma}_{sdk}\log F_{km}
  \right],
\]
normalized over $d=1,\ldots,D$. The Gamma posterior update is
\[
  \alpha_{isd}
  \leftarrow
  \alpha_{sd}^0+\sum_{m=1}^M \xi_{isdm}Y_{im},
\]
and
\[
  \beta_{sd}
  \leftarrow
  \beta_{sd}^0
  +
  \sum_{m=1}^M\sum_{k=1}^K \bar{\gamma}_{sdk}F_{km}.
\]
The second term is equal to one when each row of $F$ is normalized, but we
write the full expression to make the model dependence explicit.

For motif dimension $(s,d)$, the posterior probability of factor $k$ is updated by
\[
  \bar{\gamma}_{sdk}
  \propto
  \rho_{sdk}
  \exp\left[
    \sum_{i=1}^N \omega_{is}
    \left\{
      \sum_{m=1}^M \xi_{isdm}Y_{im}\log F_{km}
      -
      \frac{\alpha_{isd}}{\beta_{sd}}\sum_{m=1}^M F_{km}
    \right\}
  \right].
\]
This is the central posterior uncertainty object in MoS-NMF. A concentrated
vector $\bar{\gamma}_{sd1},\ldots,\bar{\gamma}_{sdK}$ means that motif dimension
$(s,d)$ has a well-identified factor. A diffuse vector means that several
factor rows are plausible explanations for the same motif dimension.

The motif responsibility for observation $i$ is updated from the component
lower bound for that observation under each motif:
\[
  \omega_{is}
  \propto
  \pi_s \exp\{\mathcal{L}_{is}\},
\]
where
\begin{align*}
  \mathcal{L}_{is}
  &=
  \sum_{d=1}^D\sum_{m=1}^M
  \xi_{isdm}Y_{im}
  \left\{
    \psi(\alpha_{isd})-\log\beta_{sd}
    +
    \sum_{k=1}^K \bar{\gamma}_{sdk}\log F_{km}
    -
    \log \xi_{isdm}
  \right\}
  \\
  &\quad
  -
  \sum_{d=1}^D\sum_{m=1}^M
  \frac{\alpha_{isd}}{\beta_{sd}}
  \sum_{k=1}^K \bar{\gamma}_{sdk}F_{km}
  -
  \sum_{d=1}^D
  \KL\{\Gam(\alpha_{isd},\beta_{sd})
  \,\|\,\Gam(\alpha_{sd}^0,\beta_{sd}^0)\}.
\end{align*}
The mixture weights are updated by
\[
  \pi_s \leftarrow \frac{1}{N}\sum_{i=1}^N \omega_{is}.
\]

\subsection{Learning the factors}

When $F$ is unknown, MoS-NMF updates $F$ using the expected allocated counts
\[
  A_{km}
  =
  \sum_{i=1}^N\sum_{s=1}^S\sum_{d=1}^D
  \omega_{is}\bar{\gamma}_{sdk}\xi_{isdm}Y_{im}.
\]
The row-normalized update is
\[
  F_{km}
  \leftarrow
  \frac{A_{km}+\epsilon}
       {\sum_{m'=1}^M(A_{km'}+\epsilon)}.
\]
In the current implementation, $F$ is initialized by MF-Poisson-SuSiE with
learned $F$, and the subsequent MoS-NMF updates alternate between motif
responsibilities, motif posteriors, Gamma loading posteriors, empirical-Bayes
prior updates, and optionally factor updates. This initialization is important
because the full problem is highly nonconvex.

\subsection{Empirical-Bayes dimension selection}

Each fitted motif starts with the same nominal dimension $D$, but the effective
dimension of motif $s$ can be smaller. MoS-NMF uses the empirical-Bayes Gamma
prior parameters to diagnose dimension activity. For each motif $s$ and
dimension $d$,
the prior mean
\[
  \frac{\alpha_{sd}^0}{\beta_{sd}^0}
\]
estimates the loading size expected for that dimension. If this value is close to
zero relative to the other dimensions in motif $s$, then dimension $d$ is effectively
inactive. This gives an ARD-style selected dimension
\[
  \widehat{D}_s
  =
  \#\left\{
    d:
    \frac{\alpha_{sd}^0/\beta_{sd}^0}
         {\sum_{d'=1}^D \alpha_{sd'}^0/\beta_{sd'}^0}
    \geq \tau
  \right\},
\]
for a threshold $\tau$.

This ARD diagnostic should be interpreted separately from the entropy of
$\bar{\gamma}_{sd1},\ldots,\bar{\gamma}_{sdK}$. The ARD prior mean asks whether
a dimension carries loading mass. The entropy of $\bar{\gamma}_{sd}$ asks whether
the factor identity of that dimension is uncertain. Both are useful: an
inactive dimension should have small loading mass, while an active but
ambiguous dimension may have large loading mass and diffuse posterior factor
identity.

\begin{algorithm}[t]
\caption{MoS-NMF variational empirical-Bayes fitting}
\small
\begin{algorithmic}[1]
\State Initialize $F$, typically using MF-Poisson-SuSiE with learned factors.
\State Initialize $\omega_{is}$ by clustering preliminary loading scores.
\State Initialize $\bar{\gamma}_{sdk}$ from cluster-level loading patterns.
\State Initialize $\alpha_{isd}$, $\beta_{sd}$, $\pi_s$, and
$(\alpha_{sd}^0,\beta_{sd}^0)$.
\For{$t=1,\ldots,T$}
  \State Update $\xi_{isdm}$ in feature blocks.
  \State Update $\alpha_{isd}$ and $\beta_{sd}$.
  \State Update $\omega_{is}$ and $\pi_s$.
  \State Update $(\alpha_{sd}^0,\beta_{sd}^0)$ by empirical Bayes.
  \State Update $\bar{\gamma}_{sdk}$ for each motif dimension.
  \If{factor learning is enabled}
    \State Update $F$ using expected allocated counts and row-normalize.
  \EndIf
  \State Monitor the variational objective or a scheduled approximation to it.
\EndFor
\State Report motif responsibilities $\omega$, motif posteriors
$\bar{\gamma}$, selected dimensions $\widehat{D}_s$, and fitted factors $F$.
\end{algorithmic}
\end{algorithm}

\section{Identifiability Theory}
\label{sec:identifiability}

This section will contain the identifiability theory. The theory should
separate at least three issues:
\begin{enumerate}
  \item permutation and scaling conventions for the NMF factorization;
  \item identifiability of motif supports given the factor matrix $F$;
  \item identifiability of the mixture of sparse motifs when motifs share
  factors or when no motif contains a single anchor factor.
\end{enumerate}
The empirical sections below are designed to stress these same issues before
the theory is complete.

\section{Experiments Needed for a JMLR-Quality Evaluation}
\label{sec:experiments}

The empirical goal is not only to show good reconstruction error. The paper
needs to show that MoS-NMF recovers meaningful sparse motifs, clusters
observations by motif, learns factors when $F$ is unknown, and quantifies
uncertainty in settings where hard sparse supports are unstable.

\subsection{Baselines}

The main baselines should be:
\begin{enumerate}
  \item Poisson NMF followed by clustering of normalized loading scores;
  \item MF-Poisson-SuSiE followed by clustering of posterior loading scores;
  \item sparse NMF baselines with tuned sparsity penalties, if available;
  \item ablations of MoS-NMF with fixed $F$, learned $F$, and different
  initialization schemes.
\end{enumerate}
The first two baselines are already implemented in the current benchmark code.

\subsection{Metrics}

The simulation tables should include:
\begin{enumerate}
  \item factor recovery, measured by cosine matching between fitted rows of
  $F$ and true rows of $F_0$;
  \item motif clustering accuracy, measured by matching $\arg\max_s\omega_{is}$
  to true groups;
  \item support recovery, measured by converting posterior motif scores to
  predicted active factors;
  \item soft clustering diagnostics, based on group-level averages of
  $\omega_{is}$ rather than only hard labels;
  \item posterior uncertainty summaries for $\bar{\gamma}_{sd}$, including
  entropy, top-two gaps, and targeted pairwise split metrics;
  \item selected motif dimensions $\widehat{D}_s$ from the ARD prior means;
  \item predictive fit, such as held-out Poisson log likelihood or deviance.
\end{enumerate}

\subsection{Core simulation settings}

The current roadmap suggests the following simulation studies.

\paragraph{Correctly specified block motifs.}
Use correct $K$ and nominal $D$, with groups occupying simple single-factor and
multi-factor supports. This verifies that the method can recover factors,
clusters, and supports in a clean setting. Current evidence shows feasibility,
but a larger seed grid is needed before claiming robust superiority over
baselines.

\paragraph{Tree/no-anchor motifs.}
Use motifs such as $(1,2,3)$, $(1,2,4)$, $(1,5,6)$, $(1,5,7)$, $(1,2)$, and
$(1,5)$, so that no group is explained by a single anchor factor. This is the
strongest current positive setting: in a small multi-seed pilot, MoS-NMF has
substantially better support recovery than the baseline pipelines.

\paragraph{Overspecified $K$.}
Generate data with $K$ true factors but fit with extra candidate factor rows,
including near-duplicates of true factors. The key diagnostic is whether
$\bar{\gamma}_{sd}$ splits posterior mass across duplicate rows when they are
statistically interchangeable. Global entropy is not sufficient; the paper
should report pair-specific posterior split metrics for the duplicated factors.

\paragraph{Correlated factors.}
Generate factor rows with controlled correlation strength. This tests whether
MoS-NMF can express targeted ambiguity between correlated factors. The current
evidence suggests that raw entropy can be misleading because weak dimensions may
become globally diffuse. The JMLR version should sweep correlation strength and
report targeted pair uncertainty, not only mean entropy.

\paragraph{Overspecified $D$.}
Fit each motif with a nominal dimension larger than the true motif dimension.
The target behavior is that extra dimensions either have small ARD prior mean
$\alpha_{sd}^0/\beta_{sd}^0$ or small posterior loading mass. The current
single-seed evidence suggests that diffuse-surplus initialization helps, while
repeated-surplus initialization duplicates active dimensions. A full paper needs
multi-seed threshold sweeps for $\widehat{D}_s$ and comparisons against the
true motif dimensions after matching fitted motifs to groups.

\subsection{Real data}

A JMLR-quality paper should include at least one real-data example where
repeated sparse loading motifs have scientific meaning. The analysis should
report the fitted factors, the motifs, motif-specific selected dimensions,
sample responsibilities, and posterior uncertainty over ambiguous motif dimensions.
The real-data example should be chosen so that the learned motifs can be
interpreted externally, not only through reconstruction quality.

\subsection{Ablations and robustness checks}

The paper should include ablations for:
\begin{enumerate}
  \item fixed versus learned $F$;
  \item initialization from Poisson NMF versus MF-Poisson-SuSiE;
  \item sensitivity to the number of motifs $S$;
  \item sensitivity to nominal $D$ and the ARD threshold $\tau$;
  \item sensitivity to duplicate-factor and correlated-factor strength;
  \item stability across random initializations.
\end{enumerate}
These experiments are needed to distinguish the proposed statistical model from
an initialization artifact.

\section{Discussion}

MoS-NMF reframes sparse NMF as a motif discovery problem. Instead of estimating
sparse loadings independently for each observation, it pools repeated sparsity
patterns through latent motifs. The posterior motif probabilities
$\bar{\gamma}_{sdk}$ provide an interpretable uncertainty summary over the
factor identities defining each motif, while the responsibilities $\omega_{is}$
cluster observations by motif. The empirical-Bayes prior means
$\alpha_{sd}^0/\beta_{sd}^0$ provide a natural ARD-style diagnostic for the
effective dimension of each motif.

The current evidence supports the promise of this approach, especially in
tree/no-anchor simulations and in overfitted-factor uncertainty diagnostics.
The main open issues are stable dimension selection, broader multi-seed
benchmarks, model-selection guidance for $K$, $S$, and $D$, and a real-data
example with scientifically interpretable motifs.

\bibliographystyle{plainnat}
\bibliography{sparsity}

\end{document}
